\scanpage{21}
\mspage{--- 49 ---}
5) Faisons maintenant l’hypothèse que $\mathrm{V}$ ne fait intervenir que les exposants 0 et 1, autrement dit que $\mathrm{V}_{\mathrm{C}}=\mathrm{V}(0)+\mathrm{V}_{\mathrm{C}}(1)$ ; posons
\[
n_{\unsure{\mathrm{o}|0}}=\dim_{\mathrm{C}}\mathrm{V}_{\unsure{\mathrm{C}|\dot{\mathrm{C}}}}(0)\ \text{et}\ n_1=\dim_{\mathrm{C}}\mathrm{V}_{\mathrm{C}}(1).
\]
(Le module $p$-adique attaché à un groupe $p$-divisible est de ce type, l’entier $n_1$ étant égal à la dimension du groupe.)

Supposons que $n_1=1$ et que $\mathrm{V}$ soit un $\mathrm{H}_{\mathrm{V}}^{\unsure{\mathrm{O}|0|\circ}}$-module simple. Alors le commutant de $\mathrm{H}_{\mathrm{V}}^{\mathrm{O}}$ dans $\operatorname{End}(\mathrm{V})$ est un corps commutatif $\mathrm{E}$ de degré fini sur $\mathbf{Q}_p$, et $\mathrm{H}_{\mathrm{V}}^{\mathrm{O}}$ est égal au groupe algébrique des $\mathrm{E}$-automorphismes de $\mathrm{V}$ ; il en résulte, conformément à ce qu’avaient conjecturé \textsc{Lubin} et \textsc{Tate}, que le groupe $\mathrm{G}_{\mathrm{V}}$ contient un sous-groupe \textit{ouvert} de $\operatorname{Aut}_{\mathrm{E}}(\mathrm{V})$.

Le cours s’est achevé par un essai de classification dans le cas $n_1>1$ (en supposant $\mathrm{V}$ absolument simple sur $\mathrm{H}_{\mathrm{V}}^{\mathrm{O}}$), basé sur le fait que $\mathrm{H}_{\mathrm{V}}^{\mathrm{O}}/_{\mathrm{C}}$ contient un tore de dimension 1 \textit{n’ayant que deux poids}.
% slanted
\margin{\handwritten{$\left.\vphantom{\begin{matrix}1\\1\end{matrix}}\right]$ voir à ce sujet le livre de T. Springer sur \unsure{les|ls} alg. de Jordan récemment paru chez Springer-Verlag}}

\begin{center}\textsc{Séminaire}\end{center}

Le séminaire a continué celui de 1965/66, relatif aux \textit{groupes $p$-divisibles}. Il a comporté neuf exposés : quatre par Michel \textsc{Raynaud} et cinq par John \textsc{Tate} (qui en avait fait dix en 1965/66).

Michel \textsc{Raynaud} a rappelé les notions préliminaires nécessaires : schémas en groupes finis, dualité de Cartier, quotients. Il a ensuite donné la définition des groupes $p$-divisibles, et montré leur interprétation comme groupes formels. Dans son dernier exposé, il a démontré le théorème fondamental, dû à \textsc{Tate}, qui affirme que, en inégale caractéristique, le passage de l’anneau des entiers au corps des fractions définit un foncteur \textit{pleinement fidèle} sur la catégorie des groupes $p$-divisibles. Il est possible que sa démonstration, différente de celle de \textsc{Tate}, puisse être adaptée au cas d’égale caractéristique.

Les exposés de John \textsc{Tate} ont eu pour thème le \textit{théorème de relèvement} des variétés abéliennes, autrement dit le fait que relever une variété abélienne revient à relever le groupe $p$-divisible correspondant (la caractéristique résiduelle étant égale à $p$). A Woods Hole, en 1964, \textsc{Tate} en avait donné une démonstration basée sur une certaine construction cohomologique. Dans le séminaire, il a proposé une autre méthode, où variété abéliennes et groupes $p$-divisibles n’interviennent que par les \textit{faisceaux} correspondants (pour des topologies d’\textsc{Artin-Grothendieck} convenables). Cette méthode, malgré des difficultés techniques qui n’ont pas été entièrement résolues, paraît plus naturelle que la précédente, et de portée plus générale.

\noindent{\footnotesize 4}
% archive sheet number: 20
